\section{Experimental Setup and Procedure}
\label{sec:procedure}

This chapter describes the experimental setup and the measurements performed.

\subsection{Experimental setup}
\label{sec:setup}

The coaxial germanium detector, shown in Figure~\ref{fig:detector-schematic}, contains a cylindrical crystal with a diameter
of \SI{45}{\milli\metre} and a length of \SI{39}{\milli\metre}.
Its outer surface forms a lithium-diffused contact, while the central bore
is coated with gold. An aluminium
cover protects the crystal. Absorption in the cover and the lithium-diffused
layer limits detection below approximately
\SIrange{40}{50}{\kilo\electronvolt}~\cite{v18-manual}.
A source holder positions the source \SI{1.5}{\centi\metre} above the
aluminium cover. The whole assembly is enclosed by by a lead shielding to reduce background radiation.

A liquid-nitrogen Dewar cools the detector, and a high-voltage supply provides
the reverse bias. The charge-sensitive amplifier converts collected charge
into a voltage signal. The analog-to-digital converter (ADC) digitises the
pulse heights, and the computer sorts them into energy-related channels
and records the histogram. The channel number depends linearly on the
deposited energy. Figure~\ref{fig:signal-chain} shows the signal
path.

\begin{figure}[!htb]
  \centering
  % Original functional signal-chain diagram.
% Requires TikZ and the arrows.meta and positioning libraries.
\begin{tikzpicture}[
  x=1cm,y=1cm,
  font=\small,
  >=Latex,
  block/.style={draw=blue!45!black,fill=blue!5,rounded corners=2pt,
    minimum height=1.05cm,align=center,inner sep=4pt},
  service/.style={draw=black!55,fill=black!3,rounded corners=2pt,
    minimum height=0.85cm,align=center,inner sep=4pt},
  signal/.style={->,line width=0.8pt},
  support/.style={->,draw=black!65,thin}
]
  \node[block,minimum width=1.65cm] (ge) at (0.9,0)
    {Ge\\detector};
  \node[block,minimum width=2.45cm] (csa) at (3.55,0)
    {Charge-sensitive\\amplifier};
  \node[block,minimum width=2.15cm] (adc) at (6.55,0)
    {ADC};
  \node[block,minimum width=2.15cm] (pc) at (9.55,0)
    {Acquisition\\computer};

  \draw[signal] (ge.east) -- (csa.west);
  \draw[signal] (csa.east) -- (adc.west);
  \draw[signal] (adc.east) -- (pc.west);

  \node[service,minimum width=2.25cm] (bias) at (0.9,1.75)
    {High-voltage supply};
  \node[service,minimum width=2.25cm] (cool) at (0.9,-1.75)
    {Liquid-nitrogen Dewar};
  \draw[support] (bias.south) -- (ge.north)
    node[midway,right,font=\footnotesize] {bias};
  \draw[support] (cool.north) -- (ge.south)
    node[midway,right,font=\footnotesize] {cooling};

\end{tikzpicture}

  \caption{Signal processing from the Ge detector to the recorded histogram,
  with the bias supply and cooling system.}
  \label{fig:signal-chain}
\end{figure}
\FloatBarrier

\subsection{Measurements}

The measurements characterise the spectrometer and apply its calibration to
activity determination and nuclide identification. Each source spectrum is
recorded for approximately one hour, while an overnight measurement without
a source provides the background spectrum. Count rates are normalised to the
recorded live time.

The $^{152}\mathrm{Eu}$ spectrum provides the energy calibration and the
energy dependence of the full-energy detection efficiency. The detector
response is further characterised with $^{137}\mathrm{Cs}$ by examining the
energy resolution, full-energy peak, and structures arising from Compton
scattering. Using the calibrated spectrometer, the activity of a
$^{133}\mathrm{Ba}$ source is then determined. Finally, an unknown sample is
investigated to identify its radioactive nuclides and determine their
activities.
